\documentclass[12pt]{article}
\usepackage{lmodern}
\usepackage{geometry}
\usepackage{enumitem}
\setlist[description]{leftmargin=1cm,labelindent=1cm}
\usepackage{amsmath}
\usepackage{amsthm}
\usepackage{amssymb}
\usepackage{mathtools}
\usepackage[portuguese]{babel}

\theoremstyle{plain}
\newtheorem{thm}{Teorema}[section]
\newtheorem{prop}[thm]{Proposição}
\newtheorem{lem}[thm]{Lema}
\newtheorem{cor}[thm]{Corolário}

\theoremstyle{definition}
\newtheorem{defn}[thm]{Definição}
\newtheorem{exmp}[thm]{Exemplo}
\newtheorem{prblm}[thm]{Problema}
\newtheorem{prcp}[thm]{Princípio}

\theoremstyle{remark}
\newtheorem*{rem}{Nota}

\geometry{
 a4paper,
 left=3cm,
 top=2cm,
 right=2cm,
 bottom=2cm
}

\begin{document}
\begin{center}
    \textbf{\large UNIVERSIDADE FEDERAL DO RIO GRANDE DO SUL} \\
    \textbf{\large INSTITUTO DE MATEMÁTICA E ESTATÍSTICA} \\
    \vspace{3cm}
    \textbf{\Large MÉTODOS APLICADOS II} \\
    \vspace{1cm}
    \textbf{\large LISTA DE RESULTADOS} \\
    \vspace{3cm}

    \textbf{LEONARDO SUAREZ}

    \vfill
    \textbf{\large Porto Alegre, RS} \\
    \textbf{\large \today}
\end{center}

\newpage
\tableofcontents
\newpage

\section{Soluções de séries de potências de equações lineares}

\begin{defn}
    Um ponto \(x=x_0\) é chamado de \emph{ponto ordinário} da equação diferencial
    \begin{equation}\label{eq:edo1}
        a_2(x)y'' + a_1(x)y' + a_0(x)y = 0,
    \end{equation}
    \noindent se ambos os coeficientes \(P(x)\) e \(Q(x)\) na forma padrão
    \begin{equation}\label{eq:edo2}
        y'' + P(x)y' + Q(x)y = 0,
    \end{equation}
\noindent são analíticos em \(x_0\). Caso contrário, \(x_0\) é chamado de \emph{ponto singular}.
\end{defn}

\begin{thm} (existência de soluções por séries de potências)
    Se \(x_0\) é ponto ordinário da equação \eqref{eq:edo1}, então podemos encontrar duas soluções linearmente independentes na forma de séries de potências centradas em \(x_0\), ou seja,
    \[
    y = \sum_{n=0}^{\infty} c_n (x - x_0)^n.
    \]
    Uma solução de série de potência converge pelo menos no intervalo \(|x - x_0| < R\), onde \(R\) é a distância do ponto \(x_0\) ao ponto singular mais próximo, ou também a \emph{cota inferior} do raio de convergência.
\end{thm}

\begin{defn}
    Uma ponto singular \(x=x_0\) é chamado de \emph{ponto singular regular} da equação \eqref{eq:edo1} se as funções \(p(x)=(x-x_0)P(x)\) e \(q(x)=(x-x_0)^2Q(x)\) são analíticas em \(x_0\). Caso contrário, \(x_0\) é chamado de \emph{ponto singular irregular}.
\end{defn}

\begin{thm} (teorema de Frobenius)
    Se \(x_0\) é ponto singular regular da equação \eqref{eq:edo1}, então existe pelo menos uma solução da forma
    \begin{equation}
        y = (x - x_0)^r \sum_{n=0}^{\infty} c_n (x - x_0)^n = \sum_{n=0}^{\infty} c_n (x - x_0)^{n+r},
    \end{equation}
    \noindent onde \(r\) é a raiz de uma equação indicial associada. A série converge pelo menos em algum intervalo \(0<x-x_0<R\).
\end{thm}

\end{document}